%% ----------------------------------------------------------------------------
%%  §1.8  Le funzioni potenza, esponenziale, logaritmo
%% ----------------------------------------------------------------------------
\section{Le funzioni potenza, esponenziale, logaritmo}
\label{sec:potenze}

\subsection{Potenze con esponente naturale e radici}

Consideriamo la \emph{funzione potenza con esponente $n\in\N$}:\index{funzione!potenza}
\begin{equation}\label{eq:potenza-n}
  f(x)=x^{n},
\end{equation}
che è definita, per ogni $x\in\R$, moltiplicando il numero $x$ per se stesso $n$ volte
(più precisamente: $x^{1}=x$ e $x^{n+1}=x^{n}\cdot x$; si pone inoltre $x^{0}=1$).

\begin{proposizione}\label{prop:potenza-crescente}
La funzione $f(x)=x^{n}$ è strettamente crescente per $x\ge0$, cioè
\begin{equation}\label{eq:potenza-crescente}
  0\le x_{1}<x_{2} \ \implica\ x_{1}^{n}<x_{2}^{n} \qquad\perogni n\in\N.
\end{equation}
\end{proposizione}

\begin{schema}
\Ipotesi{$n\in\N$; \ $0\le x_{1}<x_{2}$.}
\Tesi{$x_{1}^{n}<x_{2}^{n}$.}
\end{schema}

\begin{strategia}
Per $n=2$ il ragionamento è questo: da $x_{1}<x_{2}$, moltiplicando per $x_{1}\ge0$ si
ottiene $x_{1}^{2}\le x_{1}x_{2}$ (\cref{prop:ordine-prodotto}; uguaglianza possibile
solo se $x_{1}=0$); moltiplicando invece per $x_{2}>0$ si ottiene $x_{1}x_{2}<x_{2}^{2}$.
Mettendo insieme le due, $x_{1}^{2}<x_{2}^{2}$. Per $n=3$ si riparte da
$x_{1}^{2}<x_{2}^{2}$ e si ripete lo stesso gioco. Il «e così via» che ne segue è
proprio il contenuto del \emph{principio di induzione}, che enunceremo e useremo
rigorosamente nel \cref{sec:induzione}: qui ci accontentiamo dei primi passi.
\end{strategia}

\begin{dimostrazione}
Per $n=1$ la tesi coincide con l'ipotesi. Se $n=2$ e $0\le x_{1}<x_{2}$, allora
moltiplicando prima per $x_{1}\ge0$, poi per $x_{2}>0$, si ottiene
$x_{1}^{2}\le x_{1}x_{2}$ e $x_{1}x_{2}<x_{2}^{2}$ (la seconda disuguaglianza è stretta
perché $x_{2}\neq0$ e $x_{1}\neq x_{2}$, per la \cref{prop:semplificazione-prodotto});
cioè $x_{1}^{2}<x_{2}^{2}$.
Se $n=3$, partendo da $x_{1}^{2}<x_{2}^{2}$ si ottiene
$x_{1}^{3}=x_{1}x_{1}^{2}\le x_{1}x_{2}^{2}<x_{2}x_{2}^{2}=x_{2}^{3}$, e così via per
$n>3$ (si veda la \cref{prop:induzione-potenza} per la dimostrazione completa per
induzione).
\end{dimostrazione}

Per mezzo dell'estremo superiore dimostreremo nel prossimo capitolo
(\cref{teo:radice-n}) che a ogni $y\ge0$ corrisponde almeno un numero reale $x\ge0$ per cui
$f(x)=x^{n}=y$: cioè che $f$ è suriettiva da $[0,+\infty)$ su $[0,+\infty)$. La
condizione di stretta monotonia \eqref{eq:potenza-crescente} implica allora, per il
criterio di invertibilità (\cref{teo:monotona-invertibile}), che la funzione è
invertibile. Perciò è definita la funzione inversa di $f(x)=x^{n}$ ($x\ge0$), che si
chiama funzione \emph{radice $n$-esima}\index{radice $n$-esima} e si indica con
\begin{equation}\label{eq:radice}
  f^{-1}(x)=\sqrt[n]{x}=x^{1/n} \qquad (x\ge0).
\end{equation}
I grafici delle funzioni \eqref{eq:potenza-n}, per $x\ge0$, e \eqref{eq:radice} sono
riportati nella \cref{fig:potenze-radici}: in accordo con la \cref{fig:inversa}, sono
simmetrici rispetto alla bisettrice.

\begin{figure}[htb]
\centering
\begin{tikzpicture}
\begin{axis}[width=6cm, height=5.6cm, xmin=-0.1, xmax=1.8, ymin=-0.1, ymax=1.8,
  xtick={1}, ytick={1}, title={\footnotesize $y=x^{n}$ \ ($x\ge0$)}, title style={at={(0.4,0.92)}},
  legend entries={$n=1$,$n=2$,$n=3$}, legend pos=south east, legend cell align=left,
  legend style={font=\scriptsize, draw=none, fill=none}]
  \addplot[curva, domain=0:1.75] {x};
  \addplot[curva2, domain=0:1.32] {x^2};
  \addplot[curva3, domain=0:1.2] {x^3};
\end{axis}
\end{tikzpicture}\hfill
\begin{tikzpicture}
\begin{axis}[width=6cm, height=5.6cm, xmin=-0.1, xmax=1.8, ymin=-0.1, ymax=1.8,
  xtick={1}, ytick={1}, title={\footnotesize $y=\sqrt[n]{x}$}, title style={at={(0.4,0.92)}},
  legend entries={$n=1$,$n=2$,$n=3$}, legend pos=south east, legend cell align=left,
  legend style={font=\scriptsize, draw=none, fill=none}]
  \addplot[curva, domain=0:1.75] {x};
  \addplot[curva2, domain=0:1.75, samples=100] {sqrt(x)};
  \addplot[curva3, domain=0:1.75, samples=100] {x^(1/3)};
\end{axis}
\end{tikzpicture}
\caption{Le funzioni potenza $y=x^{n}$ (a sinistra) e le loro inverse, le radici
$y=\sqrt[n]{x}$ (a destra), per $n=1,2,3$. Tutte passano per $(0,0)$ e $(1,1)$; per
$0<x<1$ le potenze decrescono al crescere di $n$, per $x>1$ crescono.}
\label{fig:potenze-radici}
\end{figure}

\subsection{Potenze con esponente razionale e reale}

Per mezzo delle funzioni \eqref{eq:potenza-n} e \eqref{eq:radice} si può definire
l'elevazione a esponente \emph{razionale} ($m,n\in\N$, $x\in\R$, $x>0$):
\begin{equation}\label{eq:potenza-razionale}
  x^{m/n}=\sqrt[n]{x^{m}}, \qquad x^{-m/n}=\frac{1}{\sqrt[n]{x^{m}}}, \qquad x^{0}=1.
\end{equation}
Si verifica che la definizione è ben posta, cioè che il valore non dipende dalla
rappresentazione della frazione: $x^{2/4}=x^{1/2}$ (è un buon esercizio, che usa
l'iniettività della potenza $n$-esima). A questo punto è stato definito il significato
di $a^{b}$, con $a$ numero reale positivo e $b$ numero razionale. Utilizzando
l'assioma di completezza è possibile estendere la definizione di $a^{b}$ anche al caso
in cui l'esponente $b$ sia un numero reale non razionale, come vedremo nel prossimo
capitolo (\cref{sec:esponenziale-reale}): l'idea è di definire, per $a>1$, $a^{\sqrt2}$
come l'unico numero reale che separa
$\{a^{q}: q\in\Q,\ q<\sqrt2\}$ da $\{a^{q}: q\in\Q,\ q>\sqrt2\}$ --- ancora una volta
la completezza «riempie il buco». Un'altra definizione equivalente, tramite la funzione
esponenziale, sarà proposta nel capitolo sulle serie.

Elenchiamo le proprietà fondamentali delle potenze, valide per ogni coppia di basi
$a,b>0$ e ogni coppia di esponenti $s,t\in\R$ (per esponenti interi sono conseguenze
immediate della definizione; per esponenti razionali si deducono da quelle con esponente
intero; per esponenti reali saranno dimostrate nei prossimi capitoli,
\cref{teo:esponenziale} e \cref{teo:potenze-reali}):
\begin{gather}
  a^{s}\cdot a^{t}=a^{s+t}, \qquad (a^{s})^{t}=a^{s\cdot t}, \qquad (a\cdot b)^{s}=a^{s}\cdot b^{s}; \label{eq:pot1}\\
  a^{s}>0; \label{eq:pot2}\\
  a<b,\ s>0 \ \implica\ a^{s}<b^{s}; \label{eq:pot3}\\
  a<b,\ s<0 \ \implica\ a^{s}>b^{s}; \label{eq:pot4}\\
  a>1,\ s<t \ \implica\ a^{s}<a^{t}; \label{eq:pot5}\\
  a<1,\ s<t \ \implica\ a^{s}>a^{t}. \label{eq:pot6}
\end{gather}

\subsection{Funzione potenza e funzione esponenziale}

Dall'espressione $a^{b}$ derivano due diversi tipi di funzione, a seconda che si faccia
variare la base $a$ o l'esponente $b$. Nel primo caso consideriamo la \emph{funzione
potenza} $f(x)=x^{b}$, con $b\in\R$ fissato. Nel secondo caso abbiamo la \emph{funzione
esponenziale}\index{funzione!esponenziale} $f(x)=a^{x}$, con $a$ numero reale positivo
fissato.

Casi particolari della funzione potenza $f(x)=x^{b}$ sono quelli con $b=n\in\N$, oppure
$b=1/n$, già esaminati in precedenza. La funzione potenza, per $x>0$, in base alla
\eqref{eq:pot2}, è positiva. Inoltre $x^{b}$ \emph{è una funzione strettamente crescente
se $b>0$ e strettamente decrescente se $b<0$}, in base alle \eqref{eq:pot3},
\eqref{eq:pot4}. Esempi di grafici nei vari casi sono riportati nella
\cref{fig:potenze-reali}.

\begin{figure}[htb]
\centering
\begin{tikzpicture}
\begin{axis}[width=4.3cm, height=4.6cm, xmin=-0.1, xmax=2.2, ymin=-0.1, ymax=3.2,
  xtick={1}, ytick={1}, title={\footnotesize $x^{b}$, $b>1$}, title style={at={(0.4,0.92)}}]
  \addplot[curva, domain=0:1.75, samples=100] {x^2};
\end{axis}
\end{tikzpicture}\hfill
\begin{tikzpicture}
\begin{axis}[width=4.3cm, height=4.6cm, xmin=-0.1, xmax=2.2, ymin=-0.1, ymax=3.2,
  xtick={1}, ytick={1}, title={\footnotesize $x^{b}$, $0<b<1$}, title style={at={(0.45,0.92)}}]
  \addplot[curva, domain=0:2.1, samples=100] {x^0.5};
\end{axis}
\end{tikzpicture}\hfill
\begin{tikzpicture}
\begin{axis}[width=4.3cm, height=4.6cm, xmin=-0.1, xmax=2.2, ymin=-0.1, ymax=3.2,
  xtick={1}, ytick={1}, title={\footnotesize $x^{b}$, $b<0$}, title style={at={(0.5,0.92)}}]
  \addplot[curva, domain=0.56:2.1, samples=100] {x^(-2)};
  \addplot[curva2, domain=0.1:2.1, samples=100] {x^(-0.5)};
  \node[blu, font=\scriptsize, anchor=west] at (axis cs:0.75,2.2) {$b=-2$};
  \node[rame, font=\scriptsize, anchor=south] at (axis cs:1.6,0.95) {$b=-\tfrac12$};
\end{axis}
\end{tikzpicture}
\caption{La funzione potenza $y=x^{b}$ per $x>0$: crescente se $b>0$ (convessa se
$b>1$, concava se $0<b<1$), decrescente se $b<0$.}
\label{fig:potenze-reali}
\end{figure}

La \emph{funzione esponenziale} $f(x)=a^{x}$, con $a$ numero reale positivo, è una
funzione positiva, \emph{strettamente crescente se $a>1$ e strettamente decrescente se
$a<1$}, in base alle \eqref{eq:pot2}, \eqref{eq:pot5}, \eqref{eq:pot6}. Esempi di
grafici sono riportati nella \cref{fig:esponenziali}.

\begin{figure}[htb]
\centering
\begin{tikzpicture}
\begin{axis}[width=6cm, height=5cm, xmin=-2.5, xmax=2.2, ymin=-0.3, ymax=4.5,
  xtick={1}, ytick={1}, title={\footnotesize $a^{x}$, $a>1$}, title style={at={(0.3,0.9)}}]
  \addplot[curva, domain=-2.4:2.1, samples=100] {2^x};
  \addplot[tratteggio] coordinates {(1,0) (1,2) (0,2)};
\end{axis}
\end{tikzpicture}\hfill
\begin{tikzpicture}
\begin{axis}[width=6cm, height=5cm, xmin=-2.2, xmax=2.5, ymin=-0.3, ymax=4.5,
  xtick={1}, ytick={1}, title={\footnotesize $a^{x}$, $0<a<1$}, title style={at={(0.7,0.9)}}]
  \addplot[curva, domain=-2.1:2.4, samples=100] {0.5^x};
\end{axis}
\end{tikzpicture}
\caption{La funzione esponenziale $y=a^{x}$: $a=2$ (a sinistra) e $a=\tfrac12$ (a
destra). Entrambe passano per $(0,1)$, sono sempre positive, e l'una è la riflessa
dell'altra rispetto all'asse $y$, perché $(1/2)^{x}=2^{-x}$.}
\label{fig:esponenziali}
\end{figure}

Un caso notevolmente importante, e ciò diventerà chiaro nello studio delle derivate, si
ha quando la base è uguale al \emph{numero di Nepero}\index{numero di Nepero}
$e=2{,}718\,281\,8\dots$, che definiremo nel capitolo sulle successioni come limite di
$(1+1/n)^{n}$. In tal caso ovviamente si indica la funzione esponenziale con
$f(x)=e^{x}$; dato che $e>1$, la funzione $e^{x}$ è strettamente crescente.

\begin{notastorica}[Il numero $e$ e gli interessi composti]
Il numero $e$ compare per la prima volta, senza nome, nel 1618, in una tavola di
logaritmi naturali pubblicata in appendice a un'edizione dell'opera di Nepero (John
Napier, 1550--1617), probabilmente redatta da William Oughtred; fu Jakob Bernoulli, nel
1683, a incontrarlo studiando il
problema dell'\emph{interesse composto}: un capitale $C$ investito al tasso annuo del
$100\%$, con interessi capitalizzati $n$ volte l'anno, dopo un anno diventa
$C(1+1/n)^{n}$, e Bernoulli osservò che al crescere di $n$ questa quantità non cresce
all'infinito ma si avvicina a un numero compreso tra $2$ e $3$. La lettera $e$ compare
per la prima volta in una lettera di Eulero a Christian Goldbach del 1731, e sempre
Eulero dimostrò nel 1737 che $e$ è irrazionale. La trascendenza di $e$ --- il fatto che non sia radice di alcun
polinomio a coefficienti interi --- fu provata da Charles Hermite nel 1873.
\end{notastorica}

\subsection{La funzione logaritmo}

Se $a=1$, la funzione $a^{x}$ è identicamente uguale a $1$: è costante e, naturalmente,
non invertibile. Se invece $a\neq1$, la funzione esponenziale $a^{x}$ è strettamente
monotòna e quindi invertibile; si dimostra inoltre (\cref{teo:esponenziale}) che la sua
immagine è costituita da \emph{tutti} i numeri reali positivi.
La funzione inversa è quindi definita sui numeri reali positivi.

\begin{definizione}[Logaritmo]\index{logaritmo}
Sia $a>0$, $a\neq1$. La funzione inversa dell'esponenziale $x\mapsto a^{x}$, definita
per $x>0$ e a valori in $\R$, si chiama \emph{funzione logaritmo in base $a$} e si
scrive $f(x)=\log_{a}x$. Essa è dunque caratterizzata da
\begin{equation}\label{eq:log-def}
  y=\log_{a}x \ \sse\ a^{y}=x \qquad (x>0,\ y\in\R).
\end{equation}
Si suole omettere l'indicazione esplicita della base se tale base è il numero $e$:
\begin{equation}\label{eq:ln-def}
  y=\log x \ \sse\ e^{y}=x .
\end{equation}
\end{definizione}

In parole: $\log_{a}x$ è \emph{l'esponente da dare ad $a$ per ottenere $x$}. Le formule
\eqref{eq:log-def}, \eqref{eq:ln-def} sono molto usate anche nella forma seguente, che
non è altro che la \eqref{eq:inversa} scritta per questa coppia di funzioni:
\begin{equation}\label{eq:log-inversa}
\begin{aligned}
  a^{\log_{a}x}&=x, & e^{\log x}&=x & &(x>0); \\
  \log_{a}(a^{y})&=y, & \log(e^{y})&=y & &(y\in\R).
\end{aligned}
\end{equation}

\begin{proposizione}\label{prop:log-monotono}
Se la base $a$ è maggiore di $1$, il logaritmo $\log_{a}x$ è una funzione strettamente
crescente; se $0<a<1$ è strettamente decrescente.
\end{proposizione}

\begin{schema}
\Ipotesi{$a>1$ (oppure $0<a<1$); \ $0<x_{1}<x_{2}$.}
\Tesi{$\log_{a}x_{1}<\log_{a}x_{2}$ (rispettivamente $\log_{a}x_{1}>\log_{a}x_{2}$).}
\end{schema}

\begin{strategia}
È un caso particolare del \cref{teo:monotona-invertibile}: l'inversa di una funzione
strettamente crescente è strettamente crescente. Ne ripetiamo la dimostrazione per
assurdo in questo caso concreto, perché è istruttivo vedere come la monotonia
dell'esponenziale si «trasferisce» al logaritmo.
\end{strategia}

\begin{dimostrazione}
Siano $x_{1}<x_{2}$ e $y_{1}=\log_{a}x_{1}$, $y_{2}=\log_{a}x_{2}$. Se fosse
$y_{2}<y_{1}$, per la \eqref{eq:pot5} avremmo $x_{2}=a^{y_{2}}<a^{y_{1}}=x_{1}$,
contrariamente alle ipotesi. Analogamente è assurdo che $y_{1}=y_{2}$, perché avremmo
$x_{1}=x_{2}$. Quindi deve risultare $y_{1}<y_{2}$. Con dimostrazione analoga (usando
la \eqref{eq:pot6}) si verifica che il logaritmo è strettamente decrescente se la base è
minore di $1$.
\end{dimostrazione}

\begin{figure}[htb]
\centering
\begin{tikzpicture}
\begin{axis}[width=6cm, height=5cm, xmin=-0.5, xmax=4.5, ymin=-2.5, ymax=2.2,
  xtick={1}, ytick={1}, title={\footnotesize $\log_{a}x$, $a>1$}, title style={at={(0.7,0.9)}}]
  \addplot[curva, domain=0.12:4.4, samples=150] {ln(x)/ln(2)};
\end{axis}
\end{tikzpicture}\hfill
\begin{tikzpicture}
\begin{axis}[width=6cm, height=5cm, xmin=-0.5, xmax=4.5, ymin=-2.2, ymax=2.5,
  xtick={1}, ytick={1}, title={\footnotesize $\log_{a}x$, $0<a<1$}, title style={at={(0.7,0.9)}}]
  \addplot[curva, domain=0.12:4.4, samples=150] {-ln(x)/ln(2)};
\end{axis}
\end{tikzpicture}
\caption{La funzione logaritmo $y=\log_{a}x$: $a=2$ (a sinistra) e $a=\tfrac12$ (a
destra). Sono definite solo per $x>0$, passano per $(1,0)$ e sono le simmetriche
rispetto alla bisettrice delle esponenziali della \cref{fig:esponenziali}.}
\label{fig:logaritmi}
\end{figure}

\needspace{8\baselineskip}
\begin{proposizione}[Proprietà dei logaritmi]\label{prop:log-proprieta}
Sia $a>0$, $a\neq1$. Per ogni $x,x_{1},x_{2}>0$ e ogni $b\in\R$, e per ogni altra base
$c>0$, $c\neq1$, valgono:
\begin{align}
  \log_{a}(x_{1}\cdot x_{2}) &= \log_{a}x_{1}+\log_{a}x_{2}; \label{eq:log1}\\
  \log_{a}(x_{1}/x_{2})      &= \log_{a}x_{1}-\log_{a}x_{2}; \label{eq:log2}\\
  \log_{a}x^{b}              &= b\log_{a}x; \label{eq:log3}\\
  \log_{c}x                  &= \frac{\log_{a}x}{\log_{a}c} \qquad\text{(cambiamento di base).} \label{eq:log4}
\end{align}
\end{proposizione}

\begin{schema}
\Ipotesi{$a,c>0$, \ $a,c\neq1$; \ $x,x_{1},x_{2}>0$; \ $b\in\R$.}
\Tesi{Le relazioni \eqref{eq:log1}--\eqref{eq:log4}.}
\end{schema}

\begin{strategia}
Ogni proprietà del logaritmo è la «traduzione» di una proprietà delle potenze
\eqref{eq:pot1}, letta attraverso la definizione \eqref{eq:log-def}. Il metodo è sempre
lo stesso: si dà un nome agli esponenti, $y_{1}=\log_{a}x_{1}$ e $y_{2}=\log_{a}x_{2}$;
si riscrive tutto in forma esponenziale, $x_{1}=a^{y_{1}}$, $x_{2}=a^{y_{2}}$; si
applica la proprietà delle potenze; si torna ai logaritmi. Lo facciamo per la
\eqref{eq:log1} e lasciamo le altre al lettore.
\end{strategia}

\begin{dimostrazione}
Per provare la \eqref{eq:log1} poniamo
\[
  \log_{a}x_{1}=y_{1}, \qquad \log_{a}x_{2}=y_{2};
\]
in base alla definizione \eqref{eq:log-def} risulta $a^{y_{1}}=x_{1}$ e
$a^{y_{2}}=x_{2}$. Pertanto, per la prima delle \eqref{eq:pot1},
\[
  x_{1}\cdot x_{2}=a^{y_{1}}\cdot a^{y_{2}}=a^{y_{1}+y_{2}},
\]
che, di nuovo in base alla definizione \eqref{eq:log-def}, equivale a
\[
  y_{1}+y_{2}=\log_{a}(x_{1}\cdot x_{2}).
\]
Ricordando i simboli introdotti all'inizio, quest'ultima è la tesi \eqref{eq:log1}. Le
altre relazioni si dimostrano in modo analogo: per la \eqref{eq:log3} si parte da
$x^{b}=(a^{y})^{b}=a^{by}$; per la \eqref{eq:log4} si pone $y=\log_{c}x$, cioè
$x=c^{y}$, e si prende il logaritmo in base $a$ di entrambi i membri usando la
\eqref{eq:log3}: $\log_{a}x=y\log_{a}c$. Il lettore è invitato a eseguire i calcoli
esplicitamente.
\end{dimostrazione}

\begin{notastorica}[I logaritmi: trasformare prodotti in somme]
Le proprietà \eqref{eq:log1}--\eqref{eq:log3} non sono curiosità: sono il motivo per cui
i logaritmi furono inventati. All'inizio del Seicento astronomi e navigatori dovevano
eseguire a mano moltiplicazioni tra numeri di molte cifre; Nepero (1614) e, poco dopo,
Henry Briggs (1617, con la base 10) pubblicarono tavole che permettevano di sostituire
ogni prodotto con una somma: si cercano i logaritmi dei due fattori, si sommano, e si
cerca nella tavola il numero che ha quel logaritmo. Laplace disse che i logaritmi,
«abbreviando le fatiche, raddoppiarono la vita degli astronomi». Il regolo calcolatore,
usato dagli ingegneri fino agli anni Settanta del Novecento (anche per progettare
aerei e missioni spaziali), è una tavola di logaritmi incisa su due righelli
scorrevoli: sommare lunghezze proporzionali a $\log x_{1}$ e $\log x_{2}$ dà una
lunghezza proporzionale a $\log(x_{1}x_{2})$.
\end{notastorica}

\begin{python}[Esponenziali, logaritmi e le loro proprietà]
import math
# La definizione: log_a(x) è l'esponente y tale che a**y == x
a, x = 2.0, 10.0
y = math.log(x, a)                     # logaritmo in base a
print(f"log_2(10) = {y:.6f},   2**{y:.6f} = {a**y:.6f}")

# Proprietà (log1), (log3), (log4) verificate numericamente
x1, x2, b = 3.0, 7.0, 2.5
print(math.isclose(math.log(x1*x2), math.log(x1) + math.log(x2)))
print(math.isclose(math.log(x1**b), b*math.log(x1)))
print(math.isclose(math.log(x1, 10), math.log(x1)/math.log(10)))

# Il numero di Nepero come limite di (1 + 1/n)**n
for n in (1, 10, 100, 1000, 10**6):
    print(f"n = {n:>8d}   (1+1/n)^n = {(1 + 1/n)**n:.8f}")
print(f"math.e         = {math.e:.8f}")
\end{python}
\begin{risultato}
log_2(10) = 3.321928,   2**3.321928 = 10.000000
True
True
True
n =        1   (1+1/n)^n = 2.00000000
n =       10   (1+1/n)^n = 2.59374246
n =      100   (1+1/n)^n = 2.70481383
n =     1000   (1+1/n)^n = 2.71692393
n =  1000000   (1+1/n)^n = 2.71828047
math.e         = 2.71828183
\end{risultato}

\begin{esercizio}
Risolvere l'equazione $2^{x+1}=3^{x}$.
\begin{soluzione}
Prendendo il logaritmo (in qualsiasi base, per esempio naturale) di ambo i membri, che
sono positivi, e usando le \eqref{eq:log3} e \eqref{eq:log2}:
$(x+1)\log2=x\log3$, da cui $x(\log3-\log2)=\log2$ e
$x=\log2/(\log3-\log2)=\log2/\log(3/2)\approx1{,}71$. Il passaggio al logaritmo è
un'equivalenza: da $u=v$ segue $\log u=\log v$ perché $\log$ è una funzione, e da
$\log u=\log v$ segue $u=v$ \emph{perché il logaritmo è iniettivo}
(\cref{prop:log-monotono}). \emph{Scelta logica:} l'incognita compare all'esponente, e
il logaritmo, che trasforma le potenze in prodotti \eqref{eq:log3}, la «porta giù»,
riducendo il problema a un'equazione di primo grado.
\end{soluzione}
\end{esercizio}

\begin{esercizio}
Dimostrare che $\log_{a}(1/x)=-\log_{a}x$ e che $\log_{1/a}x=-\log_{a}x$.
\begin{soluzione}
Poiché $a^{0}=1$, è $\log_{a}1=0$; per la \eqref{eq:log2} con $x_{1}=1$ e $x_{2}=x$ si
ha allora $\log_{a}(1/x)=\log_{a}1-\log_{a}x=-\log_{a}x$. In particolare
$\log_{a}(1/a)=-\log_{a}a=-1$. Per la seconda identità usiamo il cambiamento di base
\eqref{eq:log4} con $c=1/a$, lecito perché $1/a>0$ e $1/a\neq1$:
$\log_{1/a}x=\log_{a}x/\log_{a}(1/a)=-\log_{a}x$. \emph{Scelta logica:} entrambe le
identità sono casi particolari delle proprietà generali \eqref{eq:log2} e
\eqref{eq:log4}: conviene riconoscere la struttura (un quoziente, un cambio di base)
invece di ripartire ogni volta dalla definizione.
\end{soluzione}
\end{esercizio}
